\section{Technical Safety Concept and System Architecture}
\label{sec:architecture}

The Synapse SbW system implements a fail-operational architecture designed to tolerate single-point hardware and software faults without compromising primary steering control. To satisfy the ASIL D requirements, both physical power paths and electrical signaling paths are fully redundant and galvanically isolated.

\subsection{Redundant System Architecture}
The system layout features two independent control channels that operate in parallel. Under nominal conditions, both channels contribute equally to active steering and cross-monitor each other. If one channel detects a critical fault, it disables itself and isolates its winding, while the standby channel immediately assumes 100\% steering control.

Figure~\ref{fig:sbw_architecture} displays the logical block diagram of the technical safety concept.

\begin{figure}[H]
    \centering
    \begin{tikzpicture}[
        box/.style={draw, rectangle, rounded corners=3pt, minimum width=2.4cm, minimum height=1.1cm, align=center, thick, fill=white, draw=primary, font=\sffamily\scriptsize},
        secbox/.style={draw, rectangle, rounded corners=3pt, minimum width=2.4cm, minimum height=1.1cm, align=center, thick, fill=white, draw=secondary, font=\sffamily\scriptsize},
        warnbox/.style={draw, rectangle, rounded corners=3pt, minimum width=2.4cm, minimum height=1.1cm, align=center, thick, fill=white, draw=warningorange, font=\sffamily\scriptsize},
        bus/.style={draw, <->, >=stealth, ultra thick, secondary},
        signal/.style={draw, ->, >=stealth, thick, primary},
        monitor/.style={draw, <->, >=stealth, thick, warningorange, dashed},
        every node/.style={font=\sffamily\scriptsize}
    ]
        % Left Column: Hand Wheel Subsystem
        \node[box] (torqA) at (0, 1.8) {\textbf{Torque/Angle A}\\\faFile~Diverse Sensor};
        \node[box] (torqB) at (0, -1.8) {\textbf{Torque/Angle B}\\\faFile~Diverse Sensor};
        
        % Middle-Left Column: HWA Controllers
        \node[secbox] (hwaA) at (3.5, 1.8) {\textbf{HWA MCU A}\\\faMicrochip~Lockstep CPU};
        \node[secbox] (hwaB) at (3.5, -1.8) {\textbf{HWA MCU B}\\\faMicrochip~Lockstep CPU};
        
        % Middle-Right Column: RWA Controllers
        \node[secbox] (rwaA) at (8.5, 1.8) {\textbf{RWA MCU A}\\\faMicrochip~Lockstep CPU};
        \node[secbox] (rwaB) at (8.5, -1.8) {\textbf{RWA MCU B}\\\faMicrochip~Lockstep CPU};
        
        % Right Column: Road Wheel Subsystem
        \node[box] (motorA) at (12, 1.8) {\textbf{PMSM Motor Winding A}\\\faMicrochip~Phase Actuator};
        \node[box] (motorB) at (12, -1.8) {\textbf{PMSM Motor Winding B}\\\faMicrochip~Phase Actuator};
        
        % Connections: Sensors to HWA
        \draw[signal] (torqA.east) -- (hwaA.west) node[midway, above] {SPI};
        \draw[signal] (torqB.east) -- (hwaB.west) node[midway, below] {SPI};
        
        % HWA Cross-Monitoring
        \draw[monitor] (hwaA.south) -- (hwaB.north) node[midway, right] {Cross-Ch. Link};
        
        % Communication Bus (FlexRay)
        \draw[bus] (hwaA.east) -- (rwaA.west) node[midway, above] {FlexRay Ch. A};
        \draw[bus] (hwaB.east) -- (rwaB.west) node[midway, below] {FlexRay Ch. B};
        
        % RWA Cross-Monitoring
        \draw[monitor] (rwaA.south) -- (rwaB.north) node[midway, left] {Cross-Ch. Link};
        
        % HWA feedback motor block (small block between HWA MCU)
        \node[warnbox] (fbmot) at (3.5, 0) {\textbf{Feedback Motor}\\\faLock~HWA Feel Actuator};
        \draw[signal] (hwaA.south) -- (fbmot.north);
        \draw[signal] (hwaB.north) -- (fbmot.south);
        
        % Connections: RWA to Motors
        \draw[signal] (rwaA.east) -- (motorA.west) node[midway, above] {PWM};
        \draw[signal] (rwaB.east) -- (motorB.west) node[midway, below] {PWM};
        
        % Watchdog monitoring the entire RWA
        \node[warnbox] (wdog) at (8.5, 0) {\textbf{Safety Monitor}\\\faLock~Window Watchdog};
        \draw[monitor] (rwaA.south) -- (wdog.north);
        \draw[monitor] (rwaB.north) -- (wdog.south);
        
    \end{tikzpicture}
    \caption{Steer-by-Wire fail-operational architecture schematic showing dual-channel signal pathways.}
    \label{fig:sbw_architecture}
\end{figure}

\subsection{Cross-Channel Monitoring and Diagnostic Concept}
To ensure instant fault containment, the microcontrollers exchange status bytes and control signals over a high-speed local serial link. The main diagnostics checked include:
\begin{itemize}
    \item \textbf{E2E Communication Safeguards:} Data packets on the FlexRay bus include sequence numbers and a 16-bit CRC to detect signal loss, corruption, or delays.
    \item \textbf{Plausibility Check of Sensor Input:} The torque values from Sensor A and Sensor B are compared. If the variance exceeds 5\% for more than 4 ms, a sensor fault state is declared.
    \item \textbf{Power Stage Diagnostics:} Gate driver diagnostic signals are checked continuously to verify proper execution of inverter switching patterns.
\end{itemize}

\subsection{Transition to Safe State}
If a diagnostic error is flagged, the system executes the safe state transition matrix described in Table~\ref{tab:safestate}.

\begin{table}[H]
    \centering
    \small
    \caption{Safe State Transition Matrix for Steer-by-Wire Control Channels.}
    \label{tab:safestate}
    \begin{tabularx}{\linewidth}{X l l p{3.8cm}}
        \toprule
        \textbf{Diagnostic Failure Trigger} & \textbf{Active State} & \textbf{Next Safe State} & \textbf{Degraded Functionality} \\
        \midrule
        Single torque sensor channel open-circuit & Nominal Dual-Ch. & Single-Channel Mode & Full steering power assist; yellow dashboard warning lit. \\
        FlexRay Channel A communication timeout & Nominal Dual-Ch. & Single-Channel Bus & Full power assist using FlexRay Channel B; steering response unchanged. \\
        RWA Inverter Phase Short-Circuit (Channel A) & Nominal Dual-Ch. & Degraded Backup Mode & Channel A isolated; Channel B drives steering with 50\% power assist. \\
        Double-channel sensor failure (Dual loss) & Nominal Dual-Ch. & Safe Shutdown State & Steering assist disabled; physical steering rack locked in neutral position. \\
        \bottomrule
    \end{tabularx}
\end{table}
